% HD-S49. GPT-6 Astra Pro /root/sixfolds; inherited effort, 21 November 2026.
\section{HD-S49: primitive monomial rank-four sources}
\label{sec:monomial-rank-four-scope}

Let \(K\subset\mathbf C\) have characteristic zero and let
\(N\subset GL_m(K)\), \(m\ge21\), be finite and absolutely
irreducible, generated by actual elements of moving rank four.
Assume that it has no nonidentity element of moving rank one,
two or three. Suppose that it preserves a complex system of
lines and that its permutation group \(P\) on those lines is
primitive. Then the diagonal kernel \(D\) is characteristic
and nonscalar, and every original line projector belongs to
\(M_m(K)\). Every actual \(K\)-linear group normalizing
\(N\) preserves this descended line system.

For the natural sixfold source of signature \((3,3)\), let
\(N\triangleleft G\subset GL_m(K)\), with \(K\) the
exact character field, and use the primitive corner of a
connected full-monodromy sphere cover and its restricted
Jacobian theta polarization. For
\[
 K=\mathbf Q(i)\quad\hbox{or}\quad
 K=\mathbf Q(\sqrt{-d}),\qquad
 d\in\{2,7,11,19,43,67,163\},                         \tag{S49.1}
\]
this branch gives ordinary determinant \([-1]\) in the
Gaussian case and cannot occur in the other seven cases.
The assertion here concerns primitive \(P\); it does not
exhaust the two-, three-, or four-dimensional complex block
alternatives of an arbitrary \(N_4\).

\paragraph{1. Exact support and the published class-eight bound.}
For an actual monomial matrix, a permutation cycle of length
\(\ell\) has fixed dimension one precisely when its cycle
holonomy is one, and zero otherwise. This follows by solving
the successive fixed-vector equations, with the final equation
being multiplication by the holonomy. Thus a moving-rank-four
generator has permutation support at most
\[
      \sum_{\ell\ge2}\ell
      \le2\sum_{\ell\ge2}(\ell-1)\le8.                \tag{S49.2}
\]
The permutation action is transitive: an orbit sum of the
specified lines would otherwise be a proper invariant
subspace. Some generator has a nonidentity permutation.
Consequently the minimum nonidentity support of \(P\),
its classical \emph{class}, is an integer from two to eight.

The elementary arguments in HD-S48, paragraph1, show that
class two gives \(S_m\) and class three gives a group
containing \(A_m\). They use actual conjugate transpositions
and three-cycles; in the latter case the support-hypergraph
induction uses the literal commutator
\([(a\,b\,c),(a\,d\,e)]=(a\,b\,d)\).
For classes four and six, Manning's original
\emph{The order of primitive groups. III}, Transactions AMS19
(1918), section38, printed p.142, gives degree at most eight
and ten. For classes five and seven, Jordan's
\emph{Recherches sur les substitutions}, JMPA series2,17
(1872), section24, printed p.367, gives degree five for class
five and degrees seven, eight or nine for class seven.
Here seven is the Mersenne case \(2^3-1\) in that exact
published statement.

The additional input is Manning,
\emph{On the Primitive Groups of Classes Six and Eight},
American Journal of Mathematics32(3) (1910),235--256,
\textbf{printed p.256, the complete eighteen-group list}:
\url{https://www.jstor.org/stable/2370074}.
Its listed degrees, read directly from the original page, are
\[
 9,9,9,10,10,10,10,11,11,12,12,13,15,15,16,16,16,16.
                                                               \tag{S49.3}
\]
Thus a primitive group of class eight has degree at most
sixteen. This use is confined to the published enumeration;
it does not claim a new proof of its classification.
All these bounds are below \(m\ge21\). Therefore
\[
                         P=A_m\quad\hbox{or}\quad S_m.
                                                               \tag{S49.4}
\]

\paragraph{2. The diagonal kernel and literal projectors.}
The argument of HD-S48, paragraphs2--3, now applies with the
same actual groups, independently of the generator rank.
We give its maps explicitly. If \(D\) were scalar, the
line stabilizer \(H_i\) has a character faithful on \(D\)
and trivial on \(H_i'\). Hence permutation projection gives
\[
                  H_i'\xrightarrow{\sim}A_{m-1}.
                                                               \tag{S49.5}
\]
An actual lift \(g\in H_i'\) of a three-cycle has order
three. Its actions on all other fixed lines are trivial:
for a fixed label \(j\ne i\) they are characters of
\(A_{m-2}\), which is perfect. On the three moved lines
its holonomy is one because \(g^3=1\). Thus
\(\operatorname{rank}(g-I)=2\), contrary to the hypothesis.
So \(D\) is nonscalar, including in a nonsplit scalar
central extension of (S49.4).

Every abelian normal subgroup of \(N\) maps trivially to
\(P\), since \(A_m\) and \(S_m\) have no nontrivial
abelian normal subgroup in these degrees. Thus \(D\) is
the unique largest abelian normal subgroup, and is
characteristic. Equality of its line characters is a
\(P\)-invariant equivalence relation. Primitivity and
nonscalarity make these characters \(\chi_j\) pairwise
distinct. Over \(E=K(\zeta_{\exp D})\) the actual
projectors are
\[
 p_j=\frac1{|D|}\sum_{t\in D}\chi_j(t^{-1})\rho(t),
 \qquad \operatorname{rank}p_j=1,\quad
 p_jp_k=\delta_{jk}p_j,\quad\sum_jp_j=I.                \tag{S49.6}
\]
They are exactly the primitive central idempotents of the
matrix image \(E[D]\). The Galois permutation of them
commutes with (S49.4). The centralizer of \(A_m\) on
these \(m\) labels is trivial, because its point stabilizer
fixes only that point. Therefore every \(p_j\) lies in
\(M_m(K)\). Characteristicness of \(D\) makes every
normalizer of \(N\) permute these same \(K\)-lines.

\paragraph{3. The actual source and polarization maps.}
Write \(U_j=p_jK^m\), \(H=\operatorname{Stab}_G(U_1)\),
choose \(v_j=\rho(t_j)v_1\), and let
\(\chi:H\to K^*\) be the actual line character. The
induced matrix rule and pushforward-function maps are
\[
\begin{aligned}
 \rho(g)v_j&=\chi(t_{g(j)}^{-1}gt_j)v_{g(j)},\\
 F_v(x)&=p_1\rho(x)v,\qquad F_v(hx)=\chi(h)F_v(x),\\
 F_{\rho(g)v}(x)&=F_v(xg),\qquad
 F\longmapsto\sum_{H\backslash G}\rho(x)^{-1}F(x).
\end{aligned}                                                   \tag{S49.7}
\]
The sheet labels \(Ht_j^{-1}\leftrightarrow t_jH\)
give the actual right-monodromy/left-deck comparison. On an
inertia cycle, the fixed-vector equation is the return-character
equation, so these maps extend to invariant stalks and
parabolic cohomology.

In the seven non-Gaussian fields in (S49.1), the finite values
of \(\chi\) are \(\pm1\). Formula (S49.7) then gives
rational traces, contradicting the exact character field.
In the Gaussian case \(\chi(H)=\mu_4\), since a smaller
image also has rational traces. Choose \(h_0\) with
\(\chi(h_0)=i\), put \(H_0=\ker\chi\), and form
\(X=\widetilde C/H_0\to Y=\widetilde C/H\).
This is the actual cyclic-four quotient of the original
cover. With the unit monomial coefficient form, choose
\(v=(1+i)v_1\), of norm two. The corner is literally
\[
 e_v=\operatorname{av}_{H_0}(1-h_0^2)/2,\qquad
 B_v=\operatorname{im}(|H|e_v:J\widetilde C\to
 J\widetilde C)^0,\qquad i_{B_v}=a_{h_0}^{-1}|_{B_v}.
                                                               \tag{S49.8}
\]
These equalities follow from the full matrix-coefficient
formula (S48.12). For the quotient \(\alpha:\widetilde
C\to X\), pullback from the cyclic primitive factor is an
actual surjective isogeny onto \(B_v\), and
\[
 \alpha_*\alpha^*=[|H_0|],\qquad
 (\alpha^*)^*\theta_{J\widetilde C}=|H_0|\theta_{JX},
 \qquad \Phi^*h_{B_v}=|H|H_V.                               \tag{S49.9}
\]
Here \(\Phi=\operatorname{PD}(j+\bar j)\),
\(j_x(g)=h(v,\rho(g)x)/2\); the two real cross terms
in Schur orthogonality give exactly \(|H|=|G|/m\).
For any originally prescribed primitive line \(Kw\), the
literal Fourier operator associated to
\(A(x)=v h(w,x)/h(w,w)\) is
\[
 b=\frac{2m}{|G|}\sum_{g\in G}
       \operatorname{Re}\operatorname{tr}(A\rho(g^{-1}))g,
 \qquad b^\dagger b=\frac2{h(w,w)}e_w.                     \tag{S49.10}
\]
Clearing its denominators by \(M\) gives the actual
\(K\)-linear isogeny from the original corner onto \(B_v\),
with multiplier \(2M^2/h(w,w)\). All positive rational
multipliers change a rank-six Hermitian determinant by a
sixth power, which is a norm.

For exactly these induced matrices, HD-S34 proves that the
parabolic lattice over \(\mathbf Z[i,1/2]\) is self-dual:
a local cycle has coinvariants \(R/(\zeta-1)\),
\(\zeta\in\mu_4\), with the explicit perfect local
pairing in the holonomy-one case, and the global relation
\((i-1)[v_1]=0\) kills coinvariants. Its negative rank-six
Hermitian determinant is a rational unit \(-2^a\).
Since \(2=N(1+i)\), (S49.9)--(S49.10) give ordinary
class \([-1]\) for the original source. The rational split
isometry is HD-S03a. Compatible flat polarization changes in
a rank-nine family use the actual common-centralizer argument
of HD-S17; no statement about an arbitrary isolated
polarization is needed.

\paragraph{4. Exact remaining scope.}
For an imprimitive permutation image, summing lines over a
nontrivial permutation block gives complex block dimension
\(b>1\). Some generator permutes these blocks, and its
moving rank is at least \(b\), whence \(b\le4\).
This proves the remaining list \(b=2,3,4\); the preceding
primitive argument makes no descent assertion for those
systems. Each requires its own actual kernel and projector
argument. No finite numerical experiment is used in the
proof above.
